% !TeX program = xelatex
% ============================================================
% Общая аннотация — документ КТ1 КОМАНДНОЙ исследовательской ВКР
% ОП «Программная инженерия» ФКН НИУ ВШЭ.
%
% Тот же жанр, что и личная аннотация (свободный ~2-стр. документ из
% полужирных лейблов, НЕ ЕСПД: без штампа/кода/ЛУ/содержания), но
% описывает исследование КОМАНДЫ В ЦЕЛОМ: общая постановка, общая цель
% и задачи, и распределение областей ответственности между соавторами.
% Частные задачи каждого участника — в его личной аннотации.
% ============================================================
\documentclass[12pt,a4paper]{extarticle}
% !TeX program = xelatex
% ============================================================
% common/annotation_preamble.tex — преамбула Аннотации (документ КТ1
% исследовательской ВКР ОП «Программная инженерия»).
%
% ЭТО НЕ ЕСПД-документ: без штампа/рамки, без кода документа, без листа
% утверждения и без содержания. Оформление по ГОСТ 7.32-2017: поля
% 30/15/20/20, Times New Roman 12 пт, интервал 1,5, обычный номер
% страницы снизу. Данные — из common/meta (\hse-макросы). Ссылки [n]
% печатаются В СТРОКУ, список источников — по ГОСТ Р 7.0.5-2008.
% ============================================================

% ── Шрифты (автогенерируются из настроек проекта) ───────────
\newcommand{\hseDefaultMono}{Courier New}
\input{../common/fonts}
\usepackage[english,russian]{babel}

% ── Геометрия по ГОСТ 7.32-2017 §6.1.1 (без места под штамп) ─
\usepackage[a4paper,left=30mm,right=15mm,top=20mm,bottom=20mm]{geometry}
\usepackage{setspace}
\onehalfspacing
\usepackage{indentfirst}
\setlength{\parindent}{1.25cm}
\setlength{\parskip}{6pt}

% ── Типографика и таблицы ───────────────────────────────────
\usepackage{microtype}
\usepackage{csquotes}
\usepackage{enumitem}
\setlist[itemize,1]{label=--}
\usepackage{graphicx}
\usepackage[table]{xcolor}
\usepackage{array}
\usepackage{tabularx}
\usepackage{booktabs}
\usepackage{float}

% ── URL и переносы ──────────────────────────────────────────
\usepackage{url}
\usepackage{xurl}
\urlstyle{same}
\usepackage{hyphenat}

% ── Цитирование: ссылки [n] В СТРОКУ (как в реальных аннотациях),
%    список источников — ГОСТ Р 7.0.5-2008. НЕ переопределяем \@cite
%    в надстрочный вид (в отличие от ЕСПД-преамбулы). ────────────
\usepackage{cite}
\makeatletter
\renewcommand\@biblabel[1]{#1.}
\makeatother

\usepackage{lastpage}
\usepackage{hyperref}
\hypersetup{
    unicode=true,
    breaklinks=true,
    colorlinks=true,
    linkcolor=black,
    urlcolor=blue,
    citecolor=blue,
    pdfborder={0 0 0}
}

% ── Данные проекта (\hse…), затем «человеческие» имена (commands) ──
% meta даёт \hseFill/\hseOptional и все \hse-макросы; commands —
% \signdateblank, \hseExecutorsBlock-зависимости (\authorsignatureline,
% \studentgroup, \studentname) для блока «Выполнил» в титуле.
\input{../common/meta}
\input{../common/commands}

% Без штампа: обычный номер страницы снизу по центру (ГОСТ 7.32 §6.3.1).
\pagestyle{plain}

\begin{document}
\sloppy
% ============================================================
% common/annotation_title.tex — титульный лист Аннотации ВКР ОП ПИ
% по официальной форме «Титул Аннотация защита темы ВКР ПИ».
%
% Лёгкий титул БЕЗ штампа/рамки/листа утверждения: УДК, две колонки
% СОГЛАСОВАНО (руководитель проекта) + УТВЕРЖДАЮ (академический
% руководитель ОП), затем «Аннотация выпускной квалификационной
% работы», тема, направление и блок «Выполнил» (в командной работе —
% состав исполнителей через \hseExecutorsBlock).
% ============================================================
\begin{titlepage}
\thispagestyle{empty}
\singlespacing
\footnotesize
\setlength{\parindent}{0pt}
\setlength{\parskip}{0pt}

\begin{center}
\textbf{ПРАВИТЕЛЬСТВО РОССИЙСКОЙ ФЕДЕРАЦИИ}\\
\textbf{ФЕДЕРАЛЬНОЕ ГОСУДАРСТВЕННОЕ АВТОНОМНОЕ}\\
\textbf{ОБРАЗОВАТЕЛЬНОЕ УЧРЕЖДЕНИЕ ВЫСШЕГО ОБРАЗОВАНИЯ}\\
\textbf{НАЦИОНАЛЬНЫЙ ИССЛЕДОВАТЕЛЬСКИЙ УНИВЕРСИТЕТ}\\
\textbf{«ВЫСШАЯ ШКОЛА ЭКОНОМИКИ»}

\vspace{0.35cm}
\hseFaculty\\
Образовательная программа «\hseSpecialization»
\end{center}

\vspace{0.25cm}
\noindent
УДК~\hseUDC

\vspace{0.4cm}
\noindent
\begin{minipage}[t]{0.46\textwidth}
\raggedright
\textbf{СОГЛАСОВАНО}\\
Руководитель проекта

\vspace{0.15cm}
\strut\hseSupervisorTitle\\[-0.05cm]
\rule{6cm}{0.4pt}\\[-0.05cm]
{\scriptsize Должность, место работы, ученая степень}

\vspace{0.12cm}
\strut\hseSupervisorName\\[-0.05cm]
\rule{6cm}{0.4pt}\\[-0.05cm]
{\scriptsize Подпись, И.О. Фамилия}

\vspace{0.1cm}
\signdateblank
\end{minipage}
\hfill
\begin{minipage}[t]{0.46\textwidth}
\raggedright
\textbf{УТВЕРЖДАЮ}\\
\strut\hseAcademicSupervisorTitle

\vspace{0.35cm}
\mbox{\rule{6cm}{0.4pt}}\\[-0.05cm]
\hseAcademicSupervisorName

\vspace{0.12cm}
«\rule{0.8cm}{0.4pt}» \rule{2.5cm}{0.4pt} \hseYear~г.
\end{minipage}

\vspace{0.95cm}
\begin{center}
\textbf{\normalsize Аннотация}

\vspace{0.12cm}
\textbf{\normalsize выпускной квалификационной работы}

\vspace{0.35cm}
на тему
\underline{\makebox[0.72\textwidth][c]{\hseProjectName}}

\vspace{0.3cm}
по направлению подготовки бакалавров \hseProgramCode\ «\hseSpecialization»
\end{center}

\vspace{0.95cm}
\noindent\hfill
\begin{minipage}[t]{0.52\textwidth}
\raggedright
% Командная работа (shared-база) — состав исполнителей; solo/личная — один.
\ifhseSharedDoc
\hseExecutorsBlock
\else
Выполнил\\
студент группы \hseGroup\\
образовательной программы\\
\hseProgramCode\ «\hseSpecialization»

\vspace{0.25cm}
\strut\hseAuthorName\\[-0.05cm]
\rule{6cm}{0.4pt}\\[-0.05cm]
{\scriptsize И.О. Фамилия}

\vspace{0.15cm}
\rule{3.6cm}{0.4pt}\quad\rule{2.2cm}{0.4pt}\\[-0.05cm]
{\scriptsize Подпись, Дата}
\fi
\end{minipage}

\vfill
\begin{center}
\textbf{\hseCity{} \hseYear}
\end{center}
\end{titlepage}
\clearpage



\begin{center}
  \textbf{\large Shared Annotation}
\end{center}

\vspace{0.4em}

\noindent\textbf{Keywords:} \hseFill{5–15 WORDS IN UPPERCASE, IN THE NOMINATIVE CASE, SEPARATED BY COMMAS, WITH NO FULL STOP AT THE END}

\medskip
\noindent\textbf{Assessment of the current state of the problem being solved.} The problem addressed by the team is~\hseFill{briefly state the shared scientific problem}. Today it is tackled with \hseFill{list the existing approaches} \cite{example-en}. Their main limitations are: \hseFill{limitation 1}; \hseFill{limitation 2}~--- which is precisely what defines the direction of the research.

\medskip
\noindent\textbf{Relevance and novelty of the research topic.} The relevance stems from \hseFill{the practical and scientific significance of the task}. The scientific novelty of the work lies in~\hseFill{what is proposed for the first time}.

\medskip
\noindent\textbf{Object of research:} \hseFill{the team's overall object of research}.

\smallskip
\noindent\textbf{Subject of research:} \hseFill{the team's overall subject of research}.

\smallskip
\noindent\textbf{Overall research objective:} \hseFill{state the overall objective of the team's work}.

\medskip
\noindent\textbf{Research tasks:}
\begin{enumerate}[label=\arabic*), leftmargin=1.3cm, topsep=2pt, itemsep=2pt]
  \item \hseFill{shared task 1};
  \item \hseFill{shared task 2};
  \item \hseFill{shared task 3};
  \item \hseFill{extend the list of tasks}.
\end{enumerate}

\medskip
\noindent\textbf{Research methods:} \hseFill{list the methods: analytical review, formal modelling, computational experiment, statistical analysis, etc.}

\medskip
\noindent\textbf{Distribution of areas of responsibility.} For~each research direction a responsible co-author is indicated; that co-author prepares and~defends the corresponding part of the work (details~--- in~their personal annotation and~thesis chapters).

\vspace{-0.2em}
\begin{table}[H]
\centering
\begin{tabularx}{\textwidth}{|>{\raggedright\arraybackslash}p{4.5cm}|>{\raggedright\arraybackslash}X|}
\hline
\textbf{Co-author} & \textbf{Research direction / task} \\
\hline
\hseFill{full name of co-author 1} & \hseFill{area of responsibility: direction, tasks, experiments} \\
\hline
\hseFill{full name of co-author 2} & \hseFill{area of responsibility: direction, tasks, experiments} \\
\hline
\end{tabularx}
\end{table}

\noindent
The co-authors of the work are~\hseAuthorsList.

\bigskip
% ── References in IEEE citation style (English track): author initials
%    first, article titles in quotes, "in Proc. …", vol./no./pp., and an
%    accessed date for online resources. ──
\noindent\textbf{References}
\vspace{-0.3em}
\renewcommand{\refname}{}
\begin{thebibliography}{99}
  \bibitem{example-en}
  A.~B. Author and C.~D. Coauthor, ``Title of the paper,'' in \emph{Proc. of the Conference}, New York, NY, USA: ACM, 2024, pp.~10--20.

  \bibitem{example-web}
  ``Resource title.'' \url{https://example.org} (accessed Jan.~1, 2026).
\end{thebibliography}


\end{document}
